% Imporditud: tools/import_eksperimendid.py
% Allikas: Eesti füüsikaolümpiaadi eksperimendi ülesannete
%   kogu (Rael Kalda, 2023), ülesanne Ü81, lahendus L81.
\setAuthor{EFO žürii}
\setRound{piirkonnavoor}
\setYear{2000}
\setNumber{G E2}
\setDifficulty{3}
\setTopic{geomeetriline optika}
\setVahendid{lääts alusel, taskulambipirn alusel, patarei, juhtmed, ekraan, mõõtejoonlaud}
\setPunktid{10}
\setAllikas{Eksperimendi kogu Ü81, lahendus lk 68}
\setLahendusseis{toores}

\prob{Läätse fookuskaugus}
Määrake kahel viisil läätse fookuskaugus. Kirjeldage oma tegevust.

\hint

\solu
Katse idee, milles on määratud kaks mõõtmise meetodit --- 1 p. a) Kumerläätse läbinud paralleelsed kiired koonduvad fookuses --- 1 p. Läätse fookuskauguse määramine --- 2 p. b) Kujutis on ümberpööratud ja sama suur ning asub kaugusel 2f , kui eseme kaugus läätsest on 2f --- 1 p. Kujutis on vähendatud, kui ese on kaugemal kui 2f --- 1 p. Kujutis on on suurendatud, kui ese on lähemal kui 2f --- 1 p. Eseme paigutamine 2f kaugusele läätsest ja kujutise leidmine --- 1 p. 2f kauguse täpsustamine --- 2 p.
\probend
